% GENERATED FILE. Edit canonical lessons, then run npm run generate:books.
% canonical-source-hash: fnv1a64:207c9f45441c6b1f
% canonical-lessons: ES-C686-cima, ES-C686-pico, ES-C686-subida, ES-C686-roca, ES-C686-escalar

\chapter{Hacia la cima}
\label{ch:es-hacia-la-cima}
\hlchaptermodality{686}

\begin{chapteropening}
\textbf{By the end of this chapter:} \emph{I can talk about climbing a mountain: the way up, the rocks, the peak and the summit.}

Complete the last lesson of chapter 686: I can talk about climbing a mountain: the way up, the rocks, the peak and the summit.
\end{chapteropening}

\section[la cima]{la cima — the top, the summit}

\label{lesson:ES-C686-cima}



\begin{quote}
Before the new one: say the Spanish for the mountain, then the Spanish for the desert.
\end{quote}

\subsection*{You'll want to know: la cima}
\textbf{la cima} — \textquotedblleft{}the top, the summit\textquotedblright{}.

\begin{quote}
Llegamos a la cima a mediodía.
\end{quote}

\begin{quote}
\emph{We reached the top at midday.}
\end{quote}

\emph{En la cima de} — at the top of.

\begin{cousinweb}
\emph{cima} comes from Latin \emph{cyma} \textquotedblleft{}a young sprout\textquotedblright{} (EtymDB), from Greek \emph{kŷma} \textquotedblleft{}a swelling, a wave\textquotedblright{}: the top of a plant, then the top of a hill.
\end{cousinweb}

\begin{grammarlens}[title={What you've built}]
One more word for telling, explaining and describing at B1.
\end{grammarlens}

\subsection*{Your turn}
\begin{itemize}
\raggedright
  \item \emph{Say it:} \emph{la cima}
  \item \emph{Say it:} \emph{la cima}, once more
  \item \emph{Say it:} a sentence of your own with \emph{la cima}
  \item \emph{Recall:} say the Spanish for the mountain, then the Spanish for the desert, then say \emph{la cima} again
\end{itemize}

\begin{tcolorbox}[breakable,colback=teal!4,colframe=teal!35!black,title={Before you move on}]
What does la cima mean? (\textquotedblleft{}The top, the summit\textquotedblright{}.) Say it once more.
\end{tcolorbox}
\section[el pico]{el pico — the peak; the beak}

\label{lesson:ES-C686-pico}



\begin{quote}
Before the new one: say the Spanish for the desert, then the Spanish for the top.
\end{quote}

\subsection*{You'll want to know: el pico}
\textbf{el pico} — \textquotedblleft{}the peak; the beak\textquotedblright{}.

\begin{quote}
Desde el pico se ve el mar.
\end{quote}

\begin{quote}
\emph{The sea can be seen from the peak.}
\end{quote}

\emph{Y pico} — and a bit: \emph{las tres y pico}.

\begin{cousinweb}
\emph{pico} is inherited from Latin \emph{beccus} \textquotedblleft{}beak\textquotedblright{} (EtymDB), which the Romans took from Gaulish.
\end{cousinweb}

\begin{grammarlens}[title={What you've built}]
One more word for telling, explaining and describing at B1.
\end{grammarlens}

\subsection*{Your turn}
\begin{itemize}
\raggedright
  \item \emph{Say it:} \emph{el pico}
  \item \emph{Say it:} \emph{el pico}, once more
  \item \emph{Say it:} a sentence of your own with \emph{el pico}
  \item \emph{Recall:} say the Spanish for the desert, then the Spanish for the top, then say \emph{el pico} again
\end{itemize}

\begin{tcolorbox}[breakable,colback=teal!4,colframe=teal!35!black,title={Before you move on}]
What does el pico mean? (\textquotedblleft{}The peak; the beak\textquotedblright{}.) Say it once more.
\end{tcolorbox}
\section[la subida]{la subida — the climb, the way up; the rise}

\label{lesson:ES-C686-subida}



\begin{quote}
Before the new one: say the Spanish for the top, then the Spanish for the peak.
\end{quote}

\subsection*{You'll want to know: la subida}
\textbf{la subida} — \textquotedblleft{}the climb, the way up; the rise\textquotedblright{}.

\begin{quote}
La subida al pico dura dos horas.
\end{quote}

\begin{quote}
\emph{The climb to the peak takes two hours.}
\end{quote}

The opposite is \emph{la bajada}, the way down.

\begin{cousinweb}
\emph{subida} is built on \emph{subir} \textquotedblleft{}to go up\textquotedblright{}, which is inherited from Latin \emph{subīre}.
\end{cousinweb}

\begin{grammarlens}[title={What you've built}]
One more word for telling, explaining and describing at B1.
\end{grammarlens}

\subsection*{Your turn}
\begin{itemize}
\raggedright
  \item \emph{Say it:} \emph{la subida}
  \item \emph{Say it:} \emph{la subida}, once more
  \item \emph{Say it:} a sentence of your own with \emph{la subida}
  \item \emph{Recall:} say the Spanish for the top, then the Spanish for the peak, then say \emph{la subida} again
\end{itemize}

\begin{tcolorbox}[breakable,colback=teal!4,colframe=teal!35!black,title={Before you move on}]
What does la subida mean? (\textquotedblleft{}The climb, the way up; the rise\textquotedblright{}.) Say it once more.
\end{tcolorbox}
\section[la roca]{la roca — the rock}

\label{lesson:ES-C686-roca}



\begin{quote}
Before the new one: say the Spanish for the peak, then the Spanish for the climb.
\end{quote}

\subsection*{You'll want to know: la roca}
\textbf{la roca} — \textquotedblleft{}the rock\textquotedblright{}.

\begin{quote}
Nos sentamos en una roca a descansar.
\end{quote}

\begin{quote}
\emph{We sat on a rock to rest.}
\end{quote}

\emph{Firme como una roca} — solid as a rock.

\begin{cousinweb}
\emph{roca} comes from Vulgar Latin \emph{rocca} (EtymDB); where that word comes from is not settled.
\end{cousinweb}

\begin{grammarlens}[title={What you've built}]
One more word for telling, explaining and describing at B1.
\end{grammarlens}

\subsection*{Your turn}
\begin{itemize}
\raggedright
  \item \emph{Say it:} \emph{la roca}
  \item \emph{Say it:} \emph{la roca}, once more
  \item \emph{Say it:} a sentence of your own with \emph{la roca}
  \item \emph{Recall:} say the Spanish for the peak, then the Spanish for the climb, then say \emph{la roca} again
\end{itemize}

\begin{tcolorbox}[breakable,colback=teal!4,colframe=teal!35!black,title={Before you move on}]
What does la roca mean? (\textquotedblleft{}The rock\textquotedblright{}.) Say it once more.
\end{tcolorbox}
\section[escalar]{escalar — to climb (a mountain, a wall)}

\label{lesson:ES-C686-escalar}



\begin{quote}
Before the new one: say the Spanish for the climb, then the Spanish for the rock.
\end{quote}

\subsection*{You'll want to know: escalar}
\textbf{escalar} — \textquotedblleft{}to climb (a mountain, a wall)\textquotedblright{}.

\begin{quote}
Mi hermana escala montañas.
\end{quote}

\begin{quote}
\emph{My sister climbs mountains.}
\end{quote}

\emph{La escalada} is climbing as a sport.

\begin{cousinweb}
\emph{escalar} is built on \emph{escala} \textquotedblleft{}ladder\textquotedblright{} (EtymDB), from Latin \emph{scāla} \textquotedblleft{}a flight of steps, a ladder\textquotedblright{}, from \emph{scandere} \textquotedblleft{}to climb\textquotedblright{}.
\end{cousinweb}

\begin{grammarlens}[title={What you've built}]
One more word for telling, explaining and describing at B1.
\end{grammarlens}

\subsection*{Your turn}
\begin{itemize}
\raggedright
  \item \emph{Say it:} \emph{escalar}
  \item \emph{Say it:} \emph{escalar}, once more
  \item \emph{Say it:} a sentence of your own with \emph{escalar}
  \item \emph{Recall:} say the Spanish for the climb, then the Spanish for the rock, then say \emph{escalar} again
\end{itemize}

\begin{tcolorbox}[breakable,colback=teal!4,colframe=teal!35!black,title={Before you move on}]
What does escalar mean? (\textquotedblleft{}To climb (a mountain, a wall)\textquotedblright{}.) Say it once more.
\end{tcolorbox}
